\documentclass[11pt]{article}
\usepackage[margin=1in]{geometry}
\usepackage{amsmath,amssymb,amsthm,mathtools}
\usepackage{booktabs}
\usepackage{graphicx}
\usepackage{microtype}
\usepackage[hidelinks]{hyperref}
\usepackage{enumitem}

\newtheorem{theorem}{Theorem}[section]
\newtheorem{proposition}[theorem]{Proposition}
\newtheorem{corollary}[theorem]{Corollary}
\newtheorem{lemma}[theorem]{Lemma}
\theoremstyle{definition}
\newtheorem{definition}[theorem]{Definition}
\theoremstyle{remark}
\newtheorem{remark}[theorem]{Remark}

\newcommand{\Q}{\mathbb Q}
\newcommand{\Z}{\mathbb Z}
\newcommand{\F}{\mathbb F}
\newcommand{\OO}{\mathcal O}
\newcommand{\tr}{\operatorname{tr}}
\newcommand{\Gal}{\operatorname{Gal}}
\newcommand{\disc}{\operatorname{disc}}
\newcommand{\Fr}{\operatorname{Fr}}
\newcommand{\Art}{\operatorname{Art}}

\title{\textbf{The Discriminant-$12$ Return}\\
\large Centered Cayley Reciprocity, Pell--Lorentz Geometry, Cyclotomic Closure, and Finite Reduction}
\author{Jeremy Ebert}
\date{2026}

\begin{document}
\maketitle

\begin{abstract}
For $g\in SL_2$ with trace $\tau$, centered part
$H=g-\tfrac{\tau}{2}I$, and Cayley transform
$K=(g-I)(g+I)^{-1}$, we prove
\[
K=\frac{2}{\tau+2}H,
\qquad
KH=\frac{\tau-2}{2}I.
\]
In the normalization $\tau=\sigma+2$, the self-reciprocity condition
$K=H^{-1}$ therefore selects the unique value $\sigma=2$, hence trace $4$
and characteristic discriminant $12$.  With the period-$[1,2]$ positive
integral normalization, the selected return is
\[
g_{12}=\begin{pmatrix}3&1\\2&1\end{pmatrix},
\qquad
\lambda=2+\sqrt3.
\]
We realize $g_{12}$ as multiplication by $\lambda$ on the ramified ideal above
$2$ in $\Q(\sqrt3)$, and its centered generator as multiplication by
$\sqrt3$.  The same unit acts on the factor-cone null coordinates by
$(x,y)\mapsto(\lambda x,\lambda^{-1}y)$, translating rapidity by
$\log\lambda$ on every shell $xy=n$; the row/column parabola tangent endpoints
are exactly its null eigenrays.  After adjoining $i$, the operator
$\mathcal Z=(H+iI)/2$ is multiplication by the primitive twelfth root
$\zeta_{12}$ on the integral cyclotomic ideal
$\mathfrak p_2\OO_{\Q(\zeta_{12})}$.  The index-$4$ inclusion
$\Z[\sqrt3,i]\subset\OO_{\Q(\zeta_{12})}$ explains the apparent half-integral
presentation.  At the ramified prime, the cyclotomic residue field is
$\F_4$; reduced multiplication by $\zeta_{12}$ has order $3$, while the
reduced Pell return is the Frobenius involution.  The $n=11$ factor point then
appears as a distinguished arithmetic specialization: its prime ideal lies in
the ramified narrow class and its cyclotomic Artin action is $T_{11}$, complex
conjugation, corresponding to the involution $J$ in the companion
representation.  No statement about zeta zeros or the Riemann hypothesis is
used in the proof.
\end{abstract}

\section{Introduction}

The purpose of this paper is to isolate a single algebraic structure that
appears simultaneously as a centered $SL_2$ return, a Pell unit, a Lorentz
boost, a cyclotomic operator, and a finite four-state symmetry.  The starting
point is deliberately elementary.  If
\[
H=g-\frac{\tr(g)}2I,
\qquad
K=(g-I)(g+I)^{-1},
\]
then $H$ and $K$ lie on the same traceless direction.  Requiring them to be
reciprocal is rigid: in the normalization $\tr(g)=\sigma+2$ it forces
$\sigma=2$.

The trace selection occurs before any integral matrix, quadratic field, Cone
geometry, or cyclotomic field is introduced.  Once trace $4$ is selected, the
period-$[1,2]$ normalization gives the primitive positive integral return
\[
g_{12}=\begin{pmatrix}3&1\\2&1\end{pmatrix},
\]
with characteristic discriminant $12$ and eigenvalues $2\pm\sqrt3$.  The
number $2+\sqrt3$ is therefore not inserted as an independent parameter: it
is the expanding eigenvalue of the selected return.

The main dependency chain is
\[
\boxed{
\begin{gathered}
K=H^{-1}
\Longleftrightarrow \sigma=2
\Longrightarrow \tau=4
\Longrightarrow \Delta=12\\
\Longrightarrow \lambda=2+\sqrt3
\Longrightarrow g_{12}=\times\lambda\text{ on }\mathfrak p_2
\Longrightarrow H=\times\sqrt3\\
\Longrightarrow \mathcal Z=\times\zeta_{12}
\text{ on }\mathfrak p_2\OO_{\Q(\zeta_{12})}
\Longrightarrow \text{finite }\F_4\text{ reduction}.
\end{gathered}}
\]
In parallel, the same Pell unit acts on the real factor cone by a Lorentz
boost.  This supplies the geometric bridge to the factor-shell coordinates
without identifying the continuous Cone with any finite quotient.

Several symmetry layers must remain distinct.  The Boolean translation group,
the ramified factor-exchange involution, the cyclotomic Galois group
$U(12)\cong V_4$, and the affine/semilinear groups of $\F_4$ are related but
are not interchangeable.  Likewise, the two four-element quotients
$\mathfrak p_2/2\mathfrak p_2$ and
$\OO_K/\mathfrak P_2\cong\F_4$ are not the same quotient ring.

\section{Main theorem}

\begin{theorem}[Discriminant-$12$ return package]
\label{thm:main}
Let $g\in SL_2(\Q)$, let $\tau=\tr(g)$, assume $g+I$ is invertible, and define
\[
H=g-\frac\tau2I,
\qquad
K=(g-I)(g+I)^{-1},
\qquad
\tau=\sigma+2.
\]
Then the following statements hold.

\begin{enumerate}[label=(\roman*)]
\item
\[
K=\frac{2}{\tau+2}H,
\qquad
KH=\frac{\tau-2}{2}I.
\]
In the hyperbolic range,
\[
K=H^{-1}\iff\sigma=2.
\]

\item At the selected trace $\tau=4$, the period-$[1,2]$ positive integral
normalization gives
\[
g_{12}=\begin{pmatrix}3&1\\2&1\end{pmatrix},
\qquad
H_{12}=g_{12}-2I
=\begin{pmatrix}1&1\\2&-1\end{pmatrix},
\]
with
\[
\chi_{g_{12}}(X)=X^2-4X+1,
\qquad
\Delta_{g_{12}}=12,
\qquad
H_{12}^2=3I.
\]

\item Let
\[
F=\Q(\sqrt3),
\qquad
\mathfrak p_2=(2,\sqrt3-1)=(1+\sqrt3).
\]
In the ordered basis $(2,\sqrt3-1)$ of $\mathfrak p_2$,
$H_{12}$ and $g_{12}$ are multiplication by $\sqrt3$ and $2+\sqrt3$,
respectively.  The primitive form
\[
q_{12}(x,y)=2x^2-2xy-y^2
\]
is preserved by $g_{12}$ and satisfies
\[
2q_{12}(x,y)=U^2-V^2,
\qquad
(U,V)=(2x-y,\sqrt3\,y).
\]

\item In factor-cone null coordinates
\[
T=\frac{x+y}{2},\qquad X=\frac{x-y}{2},\qquad Y=\sqrt{xy},
\]
so that $T^2-X^2-Y^2=0$, the same return acts by
\[
\boxed{x'=(2+\sqrt3)x,\qquad y'=(2-\sqrt3)y.}
\]
Hence every shell $xy=n$ is invariant and, with
$x=\sqrt n e^s$, $y=\sqrt n e^{-s}$,
\[
s'=s+\log(2+\sqrt3).
\]
For every factor level $u>0$, the row and column parabolas are tangent to the
fixed-$T$ circle at $T=u/2$ and at the two null eigenrays
$(T,X)=(u/2,\pm u/2)$.

\item Let
\[
K_{12}=\Q(\zeta_{12})=\Q(\sqrt3,i),
\qquad
\mathfrak P_2=\mathfrak p_2\OO_{K_{12}}.
\]
Then
\[
\mathcal Z=\frac{H_{12}+iI}{2}
\]
is multiplication by
\[
\zeta_{12}=\frac{\sqrt3+i}{2}
\]
on $\mathfrak P_2$.  In particular
\[
\Phi_{12}(\mathcal Z)=0,
\qquad
\mathcal Z^{12}=I,
\qquad
 g_{12}=-i(I+\mathcal Z)(I-\mathcal Z)^{-1}.
\]
Moreover
\[
[\OO_{K_{12}}:\Z[\sqrt3,i]]=4.
\]

\item The residue field at $\mathfrak P_2$ is
\[
\OO_{K_{12}}/\mathfrak P_2\cong\F_4.
\]
If $\omega$ is the image of $\zeta_{12}$, then
\[
\omega^2+\omega+1=0,
\qquad
\omega^3=1,
\]
so reduced multiplication by $\mathcal Z$ has order $3$.  In a suitable
additive basis of $\F_4$, the reduction of $g_{12}$ is the Frobenius
involution $z\mapsto z^2$; equivalently, in the ramified basis it is the
transvection
\[
\bar g_{12}=\begin{pmatrix}1&1\\0&1\end{pmatrix}.
\]
The order-$3$ multiplication and order-$2$ Frobenius generate
$GL_2(\F_2)\cong S_3$.
\end{enumerate}
\end{theorem}

\section{Centered Cayley reciprocity}

\begin{lemma}[Centered square]
For $g\in SL_2$ with trace $\tau$,
\[
H^2=\left(\frac{\tau^2}{4}-1\right)I.
\]
\end{lemma}

\begin{proof}
Cayley--Hamilton gives $g^2-\tau g+I=0$.  Hence
\[
\left(g-\frac\tau2I\right)^2
=g^2-\tau g+\frac{\tau^2}{4}I
=\left(\frac{\tau^2}{4}-1\right)I.
\]
\end{proof}

\begin{proposition}[Centered Cayley identity]
If $g+I$ is invertible, then
\[
K=\frac{2}{\tau+2}H,
\qquad
KH=\frac{\tau-2}{2}I.
\]
\end{proposition}

\begin{proof}
Cayley--Hamilton gives
\[
(g+I)((\tau+1)I-g)=(\tau+2)I,
\]
so
\[
(g+I)^{-1}=\frac{(\tau+1)I-g}{\tau+2}.
\]
Therefore
\[
K=\frac{(g-I)((\tau+1)I-g)}{\tau+2}
=\frac{2g-\tau I}{\tau+2}
=\frac{2}{\tau+2}H.
\]
Multiplying by $H$ and using the centered-square identity gives
\[
KH=\frac{2}{\tau+2}
\left(\frac{\tau^2}{4}-1\right)I
=\frac{\tau-2}{2}I.
\]
\end{proof}

\begin{corollary}[Unique self-reciprocal trace]
If $g$ is hyperbolic and $\tau=\sigma+2$, then
\[
K=H^{-1}\iff\sigma=2.
\]
\end{corollary}

\begin{proof}
Hyperbolicity makes $H$ invertible, so
$K=H^{-1}$ is equivalent to $KH=I$.  Thus
$(\tau-2)/2=1$, hence $\tau=4$ and $\sigma=2$.
\end{proof}

\section{The primitive integral return and the ramified ideal}

Let
\[
A_a=\begin{pmatrix}a&1\\1&0\end{pmatrix}.
\]
The period $[1,2]$ gives
\[
g_{12}=A_1A_2
=\begin{pmatrix}3&1\\2&1\end{pmatrix}.
\]
This is the distinguished primitive positive integral representative used
throughout the paper.  The Cayley condition selects the trace-$4$ hyperbolic
class; the period/positivity normalization selects this representative.

Directly,
\[
\det g_{12}=1,
\qquad
\tr g_{12}=4,
\qquad
\chi_{g_{12}}(X)=X^2-4X+1,
\]
so its eigenvalues are $2\pm\sqrt3$ and its characteristic discriminant is
$12$.

Let
\[
H_{12}=g_{12}-2I.
\]
Then
\[
H_{12}^2=3I,
\qquad
\Z[H_{12}]\cong\Z[\sqrt3]=\OO_F.
\]

\begin{proposition}[Ramified ideal realization]
Let
\[
\mathfrak p_2=(2,\sqrt3-1).
\]
In the ordered basis
\[
f_1=2,
\qquad
f_2=\sqrt3-1,
\]
multiplication by $\sqrt3$ is $H_{12}$ and multiplication by
$\lambda=2+\sqrt3$ is $g_{12}$.
\end{proposition}

\begin{proof}
One has
\[
\sqrt3 f_1=f_1+2f_2,
\qquad
\sqrt3 f_2=f_1-f_2.
\]
Thus
\[
[\times\sqrt3]_f
=\begin{pmatrix}1&1\\2&-1\end{pmatrix}=H_{12},
\]
and adding $2I$ gives
\[
[\times(2+\sqrt3)]_f=g_{12}.
\]
\end{proof}

There is a second natural basis
\[
e_1=1+\sqrt3,
\qquad
e_2=3+\sqrt3=\sqrt3 e_1.
\]
Since
\[
e_1=f_1+f_2,
\qquad
e_2=2f_1+f_2,
\]
the basis-change matrix is
\[
C_{\rm basis}=\begin{pmatrix}1&2\\1&1\end{pmatrix},
\qquad \det C_{\rm basis}=-1.
\]
Consequently
\[
C_{\rm basis}^{-1}g_{12}C_{\rm basis}
=\begin{pmatrix}2&3\\1&2\end{pmatrix},
\qquad
C_{\rm basis}^{-1}H_{12}C_{\rm basis}
=\begin{pmatrix}0&3\\1&0\end{pmatrix}.
\]
Thus the centered relation $H_{12}^2=3I$ is literally the relation
$(\sqrt3)^2=3$ on the ideal lattice.

\section{The Pell--Lorentz norm}

For
\[
\alpha=xf_1+yf_2=(2x-y)+y\sqrt3,
\]
\[
N_{F/\Q}(\alpha)
=(2x-y)^2-3y^2
=2(2x^2-2xy-y^2).
\]
Define
\[
q_{12}(x,y)=2x^2-2xy-y^2.
\]
Its discriminant is $12$, and multiplication by the norm-one unit
$2+\sqrt3$ preserves it.

With
\[
U=2x-y,
\qquad
V=\sqrt3\,y,
\]
we have
\[
\boxed{2q_{12}=U^2-V^2.}
\]
In these coordinates the return is the Lorentz boost
\[
\boxed{
\begin{pmatrix}U'\\V'\end{pmatrix}
=
\begin{pmatrix}2&\sqrt3\\\sqrt3&2\end{pmatrix}
\begin{pmatrix}U\\V\end{pmatrix}.
}
\]
Its null eigendirections have eigenvalues
\[
\lambda=2+\sqrt3,
\qquad
\lambda^{-1}=2-\sqrt3.
\]

\section{The factor cone and the real return}

Introduce the factor-cone coordinates
\[
X=\frac{x-y}{2},
\qquad
Y=\sqrt{xy},
\qquad
T=\frac{x+y}{2}.
\]
Then
\[
T^2-X^2-Y^2=0.
\]
The null coordinates of the $(T,X)$ Lorentz plane are exactly
\[
x=T+X,
\qquad
y=T-X.
\]
Therefore the discriminant-$12$ boost acts by
\[
\boxed{x'=\lambda x,\qquad y'=\lambda^{-1}y.}
\]
It follows immediately that
\[
x'y'=xy.
\]
Thus every shell $xy=n$ is invariant.

Writing
\[
x=\sqrt n e^s,
\qquad
y=\sqrt n e^{-s},
\]
gives
\[
X=\sqrt n\sinh s,
\qquad
T=\sqrt n\cosh s,
\]
and the return is the rapidity translation
\[
\boxed{s'=s+R_{12},\qquad R_{12}=\log(2+\sqrt3).}
\]
Factor exchange $(x,y)\mapsto(y,x)$ sends $X\mapsto-X$ and $s\mapsto-s$;
therefore it reverses the boost.

\subsection{Parabola--circle tangency and the null eigenrays}

For fixed factor level $u$, the row parabola in the flat $(X,Y)$ projection is
\[
Y^2=u^2+2uX,
\]
with vertex
\[
V_u^-=(-u/2,0).
\]
The reflected column parabola has vertex
\[
V_u^+=(u/2,0).
\]
A fixed-$T$ Cone slice projects to
\[
X^2+Y^2=T^2.
\]
The circle is tangent to the row/column parabola at its vertex exactly when
\[
\boxed{T=u/2.}
\]
Hence the corresponding side-view points are
\[
(T,X)=\left(\frac u2,\pm\frac u2\right),
\]
which satisfy $T^2-X^2=0$.  They are precisely the two null eigenrays of the
boost:
\[
B_{12}\binom{u/2}{u/2}
=\lambda\binom{u/2}{u/2},
\qquad
B_{12}\binom{u/2}{-u/2}
=\lambda^{-1}\binom{u/2}{-u/2}.
\]
Thus
\[
\boxed{
\text{factor-parabola tangent endpoints}
=
\text{null eigenrays of the discriminant-$12$ return}.}
\]

\section{The mod-$12$ unit shell}

The units modulo $12$ are
\[
U(12)=\{1,5,7,11\}\cong V_4.
\]
For each $r\in U(12)$, the factor pair $(r,1)$ and its mirror $(1,r)$ map to
\[
T_r=\frac{r+1}{2},
\qquad
X_r=\pm\frac{r-1}{2},
\qquad
Y_r=\sqrt r.
\]
Thus
\[
T_r^2-X_r^2=r.
\]
Explicitly,
\[
\begin{array}{c|cccc}
r&1&5&7&11\\ \hline
T_r&1&3&4&6\\
|X_r|&0&2&3&5\\
Y_r&1&\sqrt5&\sqrt7&\sqrt{11}.
\end{array}
\]
The outer pair $r=11$ lies on the anti-diagonal $x+y=12$, equivalently
$T=6$.

\begin{figure}[ht]
\centering
\includegraphics[width=\textwidth]{mod12_v4_cone_triple.png}
\caption{The mod-$12$ unit labels embedded in factor-cone coordinates.  The
geometry realizes the four labels but does not define their modular group
law.  The outer $r=11$ pair is $(11,1)$ and $(1,11)$ on $x+y=12$.}
\label{fig:mod12-unit-shell}
\end{figure}

The standard cyclotomic identification is
\[
\Gal(\Q(\zeta_{12})/\Q)\cong U(12),
\qquad
T_r:\zeta_{12}\mapsto\zeta_{12}^r.
\]
The geometric shell, the Galois $V_4$, and the Boolean translation $V_4$ are
three different structures carrying the same number of labels.

\section{Cyclotomic completion on an integral ideal}

Let
\[
K_{12}=\Q(\zeta_{12})=\Q(\sqrt3,i),
\qquad
\zeta=\zeta_{12}=\frac{\sqrt3+i}{2}.
\]
Since $H_{12}$ is multiplication by $\sqrt3$,
\[
\boxed{\mathcal Z=\frac{H_{12}+iI}{2}=\times\zeta.}
\]
This identity is initially written in a complexified rank-$2$ presentation.
The correct integral carrier is the extended ramified ideal
\[
\mathfrak P_2=\mathfrak p_2\OO_{K_{12}}.
\]
Let $\alpha=1+\sqrt3$.  Then
\[
\mathcal B_{\mathfrak P}
=\{\alpha,\alpha\zeta,\alpha\zeta^2,\alpha\zeta^3\}
\]
is a $\Z$-basis of $\mathfrak P_2$.  Because
\[
\Phi_{12}(X)=X^4-X^2+1,
\qquad
\zeta^4=\zeta^2-1,
\]
multiplication by $\zeta$ has the integral matrix
\[
\boxed{
[\times\zeta]_{\mathcal B_{\mathfrak P}}
=
\begin{pmatrix}
0&0&0&-1\\
1&0&0&0\\
0&1&0&1\\
0&0&1&0
\end{pmatrix}.}
\]
Hence the denominator in $(H+iI)/2$ is a lattice-coordinate effect, not a
failure of integrality of the cyclotomic operator.

\subsection{The index-$4$ order inclusion}

The full ring of integers is
\[
\OO_{K_{12}}=\Z[\zeta_{12}],
\]
whereas
\[
\OO_0=\Z[\sqrt3,i]
\]
is a proper suborder.  Relative to the naive basis
$\{1,\sqrt3,i,i\sqrt3\}$, the cyclotomic basis
$\{1,\zeta,\zeta^2,\zeta^3\}$ has determinant of absolute value $1/4$.
Therefore
\[
\boxed{[\OO_{K_{12}}:\OO_0]=4.}
\]
Equivalently, the discriminants satisfy
\[
\disc(\OO_0)=2304,
\qquad
\disc(\OO_{K_{12}})=144,
\]
so the index is $\sqrt{2304/144}=4$.

The additive quotient
\[
\OO_{K_{12}}/\OO_0\cong(\F_2)^2
\]
records the integral-closure gluing.  It should not be confused with the
residue field used below.

\section{Ramified reduction and the $\F_4$ model}

Reduction of the integral return modulo $2$ in the ramified basis gives
\[
\bar g_{12}
=\begin{pmatrix}1&1\\0&1\end{pmatrix},
\qquad
(x,y)\mapsto(x+y,y).
\]
The reduced norm is
\[
\bar q_{12}(x,y)=y.
\]

The integral basis change satisfies
\[
C_{\rm basis}\equiv
\Phi=\begin{pmatrix}1&0\\1&1\end{pmatrix}\pmod2.
\]
Thus the Boolean relabeling is the reduction of an integral ideal basis
change, and
\[
\Phi\bar g_{12}\Phi^{-1}
=\begin{pmatrix}0&1\\1&0\end{pmatrix}=P.
\]
In the corresponding Boolean coordinates
$(\epsilon_1,\epsilon_2)$,
\[
\bar q_{12}=\epsilon_1+\epsilon_2
=\epsilon_1\mathbin{\mathrm{xor}}\epsilon_2.
\]
Together with the Boolean translations, the factor-exchange involution gives
\[
V_4^{\rm tr}\rtimes C_2\cong D_8.
\]

This ramified four-state quotient is not a field.  The cyclotomic extension
supplies a different four-state object.  Since
\[
N_{K_{12}/\Q}(1+\sqrt3)=4,
\]
the prime ideal $\mathfrak P_2$ has residue field
\[
\boxed{\OO_{K_{12}}/\mathfrak P_2\cong\F_4.}
\]
Reducing the cyclotomic polynomial gives
\[
\Phi_{12}(X)
\equiv X^4+X^2+1
=(X^2+X+1)^2\pmod2.
\]
If $\omega=\bar\zeta$, then
\[
\omega^2+\omega+1=0,
\qquad
\omega^3=1.
\]
Hence multiplication by $\omega$ is the canonical order-$3$ action of
\[
\F_4^\times\cong C_3.
\]

In the additive basis $(1,\omega)$,
\[
[\times\omega]
=\begin{pmatrix}0&1\\1&1\end{pmatrix},
\]
while Frobenius $z\mapsto z^2$ is
\[
[\Fr]
=\begin{pmatrix}1&1\\0&1\end{pmatrix}.
\]
Thus
\[
\boxed{[\Fr]=\bar g_{12}}
\]
in this basis, and
\[
\Fr\,[\times\omega]\,\Fr^{-1}
=[\times\omega]^{-1}.
\]
Therefore
\[
\langle\times\omega,\Fr\rangle
=GL_2(\F_2)\cong S_3.
\]
With translations of the additive four-state space,
\[
\operatorname{AGL}(1,4)
=\F_4\rtimes\F_4^\times
\cong A_4,
\]
and adjoining Frobenius gives
\[
\operatorname{A\Gamma L}(1,4)
\cong S_4.
\]

\begin{remark}[Two four-state quotients]
The ramified quotient
\[
\mathfrak p_2/2\mathfrak p_2
\]
and the residue field
\[
\OO_{K_{12}}/\mathfrak P_2\cong\F_4
\]
have the same underlying dimension over $\F_2$ but different ring
structures.  The former is nonreduced; the latter is a field.  The exact
bridge used here is the compatible linear action: the ramified transvection
matrix becomes Frobenius in the $\F_4$ model.
\end{remark}

\section{The distinguished $n=11$ specialization}

The preceding construction does not require choosing $n=11$.  The number
$11$ enters afterward as a distinguished arithmetic specialization of the
already selected discriminant-$12$ structure.

Geometrically,
\[
(11,1),(1,11)
\]
lie on the factor shell $xy=11$ and on the anti-diagonal
\[
x+y=12.
\]
Their Cone coordinates are
\[
\boxed{(X,Y,T)=(\pm5,\sqrt{11},6).}
\]
They are ordinary points of the invariant hyperbolic shell, not fixed points
of the Pell return.

Arithmetically,
\[
N_{F/\Q}(1+2\sqrt3)=-11.
\]
Let
\[
\mathfrak p_{11}=(1+2\sqrt3).
\]
Since
\[
(1+\sqrt3)(1+2\sqrt3)=7+3\sqrt3
\]
is totally positive, the two ideals determine the same narrow ideal class:
\[
\boxed{[\mathfrak p_{11}]=[\mathfrak p_2].}
\]
For the narrow Hilbert class field
\[
K_{12}=\Q(\zeta_{12})
\]
of $F=\Q(\sqrt3)$, the Artin action of this class is the cyclotomic
involution labeled by $11$:
\[
\boxed{\Art_{K_{12}/F}(\mathfrak p_{11})=T_{11}.}
\]
Because
\[
11\equiv-1\pmod{12},
\]
\[
T_{11}(\zeta_{12})
=\zeta_{12}^{11}
=\zeta_{12}^{-1}
=\overline{\zeta_{12}}.
\]
Thus $T_{11}$ is complex conjugation in the cyclotomic realization.

In the companion operator representation, this involution corresponds to the
operator $J$.  We therefore record the representation correspondence
\[
\boxed{T_{11}\longleftrightarrow J.}
\]
This notation does not identify a field automorphism literally with a matrix
or operator in another category; it states that the two realize the same
involution under the established representation.

The full specialization chain is therefore
\[
\boxed{
11
\longrightarrow
\mathfrak p_{11}
\longrightarrow
[\mathfrak p_{11}]=[\mathfrak p_2]
\longrightarrow
T_{11}
\longleftrightarrow
J.}
\]
The role of $11$ is consequently exemplary rather than foundational: the
discriminant-$12$ return predicts the setting in which its geometric,
narrow-class, and cyclotomic properties coincide.

\section{Symmetry layers and scope}

The construction contains several finite groups that should not be conflated:
\[
\begin{array}{c|c}
U(12)\cong V_4 & \text{cyclotomic Galois labels},\\
V_4^{\rm tr} & \text{Boolean translations},\\
V_4^{\rm tr}\rtimes C_2\cong D_8 & \text{translations plus factor exchange},\\
\F_4^\times\cong C_3 & \text{nonzero residue-field multiplication},\\
\operatorname{AGL}(1,4)\cong A_4 & \text{affine field symmetry},\\
\operatorname{A\Gamma L}(1,4)\cong S_4 & \text{affine semilinear symmetry}.
\end{array}
\]
Their coexistence is part of the discriminant-$12$ package, but equality of
orders or cardinalities does not identify the groups or their actions.

A separate analytic companion may realize the same Pell unit at a certified
operator crossing.  Such a realization is not used here.  In particular,
this paper proves no statement about the location of zeta zeros and makes no
claim concerning the Riemann hypothesis.

\section{Conclusion}

Centered Cayley reciprocity selects trace $4$ before arithmetic data are
introduced.  The positive integral normalization then exposes the Pell unit
$2+\sqrt3$, whose action can be read in four compatible ways:
\[
\boxed{
\begin{array}{ccl}
\text{integral return} &:& g_{12}=\times(2+\sqrt3)\text{ on }\mathfrak p_2,\\
\text{centered generator} &:& H_{12}=\times\sqrt3,\\
\text{real geometry} &:& (x,y)\mapsto(\lambda x,\lambda^{-1}y),\\
\text{cyclotomic closure} &:& \mathcal Z=\times\zeta_{12}\text{ on }\mathfrak P_2.
\end{array}}
\]
At the ramified prime, the same structure separates into an order-$3$
cyclotomic multiplication and an order-$2$ Frobenius return on $\F_4$.  The
$n=11$ shell then supplies a distinguished arithmetic example in which the
outer $x+y=12$ factor point, the ramified narrow class, cyclotomic complex
conjugation $T_{11}$, and the companion involution $J$ meet without being used
to define the theory.

\end{document}
